\documentclass[a4paper]{article}
\usepackage[pdftex]{graphics}
\usepackage[utf8]{inputenc}
\usepackage{enumerate}
\usepackage{amssymb}
\usepackage{icomma}
\usepackage{tikz}
\usepackage{href-ul}
\hypersetup{
	colorlinks=true,
	linkcolor=blue,
	urlcolor=blue}
\usepackage{geometry}
\geometry{a4paper, top=15mm, left=15mm, right=15mm, bottom=15mm,
	headsep=10mm, footskip=12mm}

\begin{document}
\begin{tikzpicture}(10,3)
	\draw[ultra thick](2,0) --(10,0) -- (10,1.3) --(2,1.3) -- (2,0);
	\draw[fill=black](2,0)-- (0,.6) -- (2,.6) -- (2,0);
	\node at (6,.6) {\large Lösungsblatt von \href{https://www.okuyakl.de}{www.okuyakl.de}};
\end{tikzpicture}
\vspace{0.5 cm}

\noindent{\bf Aufgabe 1.}\\
\begin{minipage}{0.5\textwidth}
	\includegraphics[width=7 cm]{gbrf049}
\end{minipage}
\hfill
\begin{minipage}{0.5\textwidth}
	$$
	\begin{array}{rcll}
		{1 \over x-2} - 1 &=& 0 & | +1\\
		{1 \over x-2} &=& 1     & | \cdot (x-2)\\
		1 &=&(x-2 \\|+2\\
		3 &=&  x  
	\end{array}
	$$
	Die Nullstelle liegt bei $N(3|0)$
\end{minipage}


\noindent {\bf Aufgabe 2.}\\
Der Faktor $(x-2)$ kommt auch im Nenner $(2x-4)=2(x-2)$ vor. Der Hauptnenner ist also $2(x-4)$.
Wir erweitern also lediglich den zweiten Bruch mit $2$ und erhalten zwei gleichnamige Brüche, die wir wie folgt addieren:
$${x \over 2(x-2)} + {2 \over 2(x-2)}= {x+2 \over 2(x-2)} $$
Dieses Ergebnis erhalten wir ebenfalls, wenn wir den Bruch aus der Aufgabenstellung kürzen:
$${x^2-4\over (2x-4)(x-2)}= {(x+2)(x-2)\over 2(x-2)^2}= {x+2 \over 2(x-2)}$$  
\vspace{0.5 cm}

\noindent{\bf Aufgabe 3.}\\

\renewcommand{\arraystretch}{2}
\begin{tabular}{p{3 cm}p{3 cm}p{3 cm}p{3 cm}}
a) $x=-2$ & b) $y=4$ & c) $x=1$ & d) $x=-{1\over5}$ \\
e) $x=2$ & f) $x=3$ & g) $x=5$ & h) $x=3$\\
i) $x=10$ & j) $x=-18$ & k) $x=-{1\over2}$ &
\end{tabular}
\vspace{0.5 cm}

\noindent {\bf Aufgabe 3. l)}\\
Wir faktorisieren beide Nenner:
$$
\renewcommand{\arraystretch}{2}
\begin{array}{rcll}
{1 \over 2(3x-1)} &=& 1 + {x \over 3(3x-1)} &| \cdot 6(3x-1)\\
3 &=& 6(3x-1) + 2x \\
3 &=& 18x-6 +2x &|+6 \\
9 &=& 20x &|:20 \\
0,45 &=& x
\end{array}
$$
\vspace{0.5 cm}

\noindent {\bf Aufgabe 4.}\\

\renewcommand{\arraystretch}{2}
\begin{tabular}{p{3.5 cm}p{3.5 cm}p{3.5 cm}p{3.5 cm}}
a) $1$	&b) $-{1\over 2x+2}$ &c) ${2\over(x+1)(x+3)}$ &d) ${2\over x+5}$ \\
e) $x+1$ & f) ${2x+1 \over 1-x}$ &g) ${1\over 3-x}$ &h) $3x-24$ \\
i) ${4 \over x+1}$ & j) ${-x-2 \over 2} $ 
\end{tabular}
\vspace{0.5 cm}

\noindent {\bf Aufgabe 4. k)}\\
$${5x \over x-2} \cdot {-2(x-2) \over x(x+2)}=-{10 \over x+2}$$
\vspace{0.5 cm}

\noindent {\bf Aufgabe 4. l)}\\
$$\left({1 \over 2x}+{2 \over x}\right): {3 \over x-1}= {1+4 \over 2x} \cdot {x-1 \over 3} = {5(x-1) \over 6x}$$
\vspace{0.5 cm}

\noindent{\bf Aufgabe 5.}\\
\renewcommand{\arraystretch}{2}
\begin{tabular}{p{3 cm}p{3 cm}p{3 cm}p{3 cm}}
a) $x$ &b) ${2\over x^2} $ &c) ${1\over 2x} $ &d) ${x^2 \over 49}$ \\
e) ${1\over 4x+4} $ & f) ${2\over 7x} $ &g) $1 $ &h) ${t^3 \over 4} $ \\
i) $x-1 $ & j) ${3\over 2}$ & k) ${1\over 4y-2} $ & l) ${1\over2}(x+1) $ \\
\end{tabular}

\begin{center}
	\includegraphics[width=7 cm]{../../viecher/endcomic.pdf}
	
	Hier geht es zurück zum \href{https://www.okuyakl.de/math/m8brtL049/aa049.pdf}{Aufgabenblatt}
\end{center}

\end{document}

